\section{Proofs of Theoretical Results}
\label{sec:proofs}

\begin{proof}[Proof of \cref{prop:offline}]
\textbf{(i).}\;
Under $\mathcal{D}_t(c,x_t)=\mathcal{D}_t$,
$G_t(K_t)$ depends only on $(t,K_t)$, so
$\sum_t G_t(K_t)$ is the same function of $\{K_t\}$ for every $(c,x)$.
Any $\{K_t^*\}$ solving~\eqref{eq:alloc} is therefore simultaneously
optimal for every prompt.

\textbf{(ii).}\;
Treating $K$ as continuous,
under the location-scale assumption~\eqref{eq:gain-factor},
$G_t(K)=\sigma_t\,a_K$ is strictly concave in $K$
(since $a_K$ is strictly concave; see David and Nagaraja, 2003,
Ch.~4 for the Gaussian case; the integer solution is obtained by
rounding and loses at most one unit of gain per timestep).
The Lagrangian of~\eqref{eq:alloc} is
$\mathcal{L}=\sum_t G_t(K_t)-\lambda\bigl(\sum_t K_t - B\bigr)
+\sum_t \mu_t(K_t-1)$,
giving the first-order and complementary-slackness conditions
\[
  G_t'(K_t^*)=\lambda^*-\mu_t^*,
  \qquad
  \mu_t^*\ge 0,
  \qquad
  \mu_t^*(K_t^*-1)=0.
\]
If $K_t^*>1$ then $\mu_t^*=0$ and $G_t'(K_t^*)=\lambda^*$.
Substituting $G_t'(K)=\sigma_t\,a_K'$:
\[
  \sigma_t\,a_{K_t^*}' = \lambda^*
  \quad\Longrightarrow\quad
  a_{K_t^*}' = \frac{\lambda^*}{\sigma_t}.
\]
Because the location-scale assumption makes $a_K$ the \emph{same}
function for every~$t$, the right-hand side pins $K_t^*$ on a single
strictly decreasing curve $a_K'$.
Hence $\sigma_{t_1}>\sigma_{t_2}$
implies $a_{K_{t_1}^*}' < a_{K_{t_2}^*}'$, hence $K_{t_1}^*>K_{t_2}^*$.
(If the shape of $\mathcal{D}_t$ varied with~$t$, each timestep would
have its own curve $a_{K,t}'$ and this monotonicity would not follow
from $\sigma_{t_1}>\sigma_{t_2}$ alone.)

\textbf{(iii).}\;
At $K_t\equiv B/T$, the marginal gains are
$G_t'(B/T)=\sigma_t\,a_{B/T}'$.
If $\sigma_{t_1}\neq\sigma_{t_2}$ for some pair, these are not equal,
violating the KKT condition.
Consider the perturbation $K_{t_1}=B/T+\epsilon$,
$K_{t_2}=B/T-\epsilon$ with all other $K_t$ unchanged.
By Taylor expansion,
\begin{align*}
  &\;\bigl[G_{t_1}(B/T+\epsilon)+G_{t_2}(B/T-\epsilon)\bigr]
   -\bigl[G_{t_1}(B/T)+G_{t_2}(B/T)\bigr] \\
  &= \epsilon\bigl[G_{t_1}'(B/T)-G_{t_2}'(B/T)\bigr]+O(\epsilon^2) \\
  &= \epsilon\,(\sigma_{t_1}-\sigma_{t_2})\,a_{B/T}'+O(\epsilon^2)
   \;>\;0
\end{align*}
for small $\epsilon>0$.
The common factor $a_{B/T}'$ can be pulled out because $a_K$ is
$t$-independent (location-scale); combined with
$\sigma_{t_1}-\sigma_{t_2}>0$ and $a_{B/T}'>0$, the perturbation
is strictly positive.
The uniform allocation is therefore not a stationary point.
Since the objective is concave over the convex constraint set
$\{\sum_t K_t=B,\;K_t\ge 1\}$, any non-stationary feasible point
is strictly suboptimal, so the global maximizer $\{K_t^*\}$
must satisfy $\sum_t G_t(K_t^*)>\sum_t G_t(B/T)$ whenever
$\{\sigma_t\}$ are not all equal.
\end{proof}

\begin{proof}[Proof of \cref{prop:online}]
\textbf{(i).}\;
For fixed $\{K_t\}$ with $\sum_t K_t=B$ and $K_t\ge 1$, the map
$\boldsymbol{\sigma}\mapsto\sum_t\sigma_t\,a_{K_t}$ is linear.
Hence
\[
  V^*(\boldsymbol{\sigma})
  \;=\;
  \max_{\{K_t\}:\,\sum_t K_t=B,\;K_t\ge 1}
  \sum_t\sigma_t\,a_{K_t}
  \;=\;
  \sup_{\{K_t\}\in\mathcal{K}}\;
  \langle\boldsymbol{a}_K,\,\boldsymbol{\sigma}\rangle
\]
is the pointwise supremum of linear functions, therefore convex
(Boyd and Vandenberghe, 2004, \S3.2.3).
By Jensen's inequality,
\[
  \mathbb{E}_{\boldsymbol{\sigma}}
  \bigl[V^*(\boldsymbol{\sigma})\bigr]
  \;\ge\;
  V^*\!\bigl(\mathbb{E}[\boldsymbol{\sigma}]\bigr)
  \;=\;
  V^*(\bar{\boldsymbol{\sigma}}).
\]
Since $\mathcal{K}$ is finite (integer allocations), $V^*$ is the
maximum of finitely many linear functions, hence convex and
piecewise-linear.
Equality holds iff the optimizer
$\{K_t^*(\boldsymbol{\sigma})\}$ is constant a.s., i.e.,
$\boldsymbol{\sigma}(c,x)=\bar{\boldsymbol{\sigma}}$ a.s.

\textbf{(ii).}\;
Under the location-scale model~\eqref{eq:gain-factor},
$G_t(K)=\sigma_t\,a_K$ yields
\[
  G_t'(K)=\sigma_t\cdot a_K'.
\]
Because $a_K$ is $t$-independent, the marginal gain separates into
a timestep-specific factor~$\sigma_t$ and a universal exhaustion
curve~$a_K'$; this separation is what allows the two thresholds below
to monitor each factor independently.
The optimal stopping condition is $G_t'(K)\le\lambda^*$, where
$\lambda^*$ is the shadow price of the budget constraint.
For each completed step $s$, define
$\bar{g}_s=G_s(\hat{K}_s)/(\hat{K}_s - 1)$.
Since $G_s$ is concave with $G_s(1)=0$
(a single candidate yields no gain), the chord from
$(1,0)$ to $(\hat{K}_s, G_s(\hat{K}_s))$ lies above the tangent:
\[
  \bar{g}_s
  =\frac{G_s(\hat{K}_s)-G_s(1)}{\hat{K}_s-1}
  \ge G_s'(\hat{K}_s)
  \qquad\text{(chord}\;\ge\;\text{tangent)}.
\]
Averaging:
\[
  \bar{g}
  =\frac{1}{|\mathcal{H}_g|}\sum_s\bar{g}_s
  \;\ge\;
  \frac{1}{|\mathcal{H}_g|}\sum_s G_s'(\hat{K}_s).
\]
If the realized allocations $\{\hat{K}_s\}$ coincide with the optimal
solution of~\eqref{eq:alloc}, then
$G_s'(\hat{K}_s)=\lambda^*$ for every active $s$ (by Prop.~\ref{prop:offline}(ii)),
giving $\bar{g}\ge\lambda^*$.
In practice $\{\hat{K}_s\}$ are determined online and need not be
exactly optimal, so $\beta_g\bar{g}$ ($\beta_g\in(0,1)$) serves as a
conservative, adaptive estimate of $\lambda^*$.
The gain threshold $\widetilde{g}_t^{(j)}<\beta_g\bar{g}$ tests the
product $\sigma_t\,a_K'\le\lambda^*$; the variance threshold
$\hat{\sigma}_t^{(j)}<\beta_\sigma\bar{\sigma}$ independently tests
$\sigma_t\le\bar{\sigma}$.
Requiring both thresholds to be met before stopping avoids
premature termination at high-value timesteps under estimation noise.
If only the gain threshold is checked, a timestep with high
$\sigma_t$ but temporarily low $a_K'$ (exhaustion from a few
unlucky draws) would be abandoned even though a single additional
candidate could yield a large jump.
Conversely, if only the variance threshold is checked, a low-$\sigma_t$
timestep would be abandoned even when its $a_K'$ is still large and
the search has not yet diminished.
The conjunction ensures that stopping occurs only when both the
timestep's intrinsic sensitivity and its remaining search potential
are genuinely low, preventing the loss of high-value search
opportunities.
\end{proof}
